\section{실험 환경 및 평가 방법}

\subsection{실험 환경}

세 오케스트레이션 구조는 동일한 멀티 에이전트 시스템 위에 구현되었으며,
에이전트 정의·LLM 설정·판단 로직을 공유한다. 실험 환경은
표~\ref{tab:env}와 같다.

\begin{table}[t]
\centering
\caption{실험 환경}
\label{tab:env}
\begin{tabular}{ll}
\toprule
\textbf{항목} & \textbf{내용} \\
\midrule
LLM        & Gemma (gemma4:31b, 로컬 Ollama) \\
Framework  & Google ADK 기반 자체 시스템 \\
Task       & 문서(챕터) 생성 \\
Agent      & Planner, Writer, Reviewer, Reviser \\
반복 횟수  & 작업·구조별 20회 \\
\bottomrule
\end{tabular}
\end{table}

작업 특성에 따른 차이를 관찰하기 위해 세 가지 유형의 입력 작업을 사용한다:
(1) \textbf{정형 반복형}(예: 정형화된 기술 보고서), (2) \textbf{창의형}
(예: 자유로운 서술), (3) \textbf{혼합형}(정형 구조 + 비정형 서술). 각
작업 $\times$ 각 구조 조합을 20회 반복하여 일관성 지표를 산출한다.

\subsection{비교 항목}

평가 지표는 표~\ref{tab:metric}과 같다. 품질은 동일한 평가 프롬프트를
사용하는 LLM-as-Judge~\cite{zheng2023judge} 방식으로 0$\sim$100점으로
채점한다. 일관성은 동일 입력 반복 실행 시 점수 및 산출물 길이의 표준편차로
측정한다.

\begin{table}[t]
\centering
\caption{비교 지표}
\label{tab:metric}
\begin{tabular}{ll}
\toprule
\textbf{Metric} & \textbf{설명} \\
\midrule
Execution Time & 생성 소요 시간 (초) \\
Token Usage    & 입력+출력 토큰 합계 \\
Quality        & LLM-as-Judge 점수 (0--100) \\
Consistency    & 반복 결과의 표준편차 \\
Retry Count    & 재작성/재검토 호출 횟수 \\
\bottomrule
\end{tabular}
\end{table}
